\tzpanchor{0007}
\tzpentryheader{0007}{Temporal Displacement Factor}

\begingroup\small
\begingroup\scriptsize
\setlength{\tabcolsep}{4pt}
\renewcommand{\arraystretch}{0.92}
\noindent\begin{tabularx}{\linewidth}{@{}p{0.24\linewidth}p{0.36\linewidth}X@{}}
\textbf{UUID} & \multicolumn{2}{l@{}}{\tzpcode{79af348c-4fcf-5f82-b585-f48f4326b3c3}}\\
\textbf{PiID} & \tzpcode{5358979323846} & \textbf{TZPi}\quad \tzpcode{9}\quad \textbf{PiPE}\quad \tzpcode{14}\\
\textbf{RTEID} & \textbf{Poloidal $\theta$}\quad \tzpcode{27.1826 rad} & \textbf{Toroidal $\phi$}\quad \tzpcode{00.0448 rad}\\
\textbf{TZPID} & \tzpcode{ID0007} & \textbf{Node Dependencies}\quad \tzpidref{0001}\\
\textbf{Registry} & \tzpcode{registry\_id\_0007} & \textbf{Scale}\quad temporal\\
\textbf{LOC Classification} & QC173.59 relativity and time & QA641 metric geometry\\
\textbf{arXiv Classification} & Primary: gr-qc \quad Cross-list: math-ph, hep-th & \\
\textbf{Primary Support} & Local-to-Global Time Ratio & Approach-to-TZP Limit\\
\textbf{Closure Support} & Normalized Temporal Operator & \\
\end{tabularx}\par
\endgroup
\endgroup

\tzpsection{Canonical Statement}
The temporal displacement factor $\AlphaT$ is defined by $\AlphaT = \frac{dt'}{dt}$, where $t'$ is the local time measured along a geodesic approaching $p$ and $t$ is the global time parameter.

\tzpsection{Canonical Equation}
\[
\alpha_{T}
=
\frac{dt'}{dt}
\]
\[
\alpha_{T}^{\mathrm{TZP}}
=
\lim_{\gamma \to p} \frac{dt'}{dt}
\]

\begin{minipage}[t]{0.49\linewidth}
\tzpsection{Canonical Symbol Key}
\begingroup
\small
\raggedright
\textbf{Source equation symbols.}\par
\begin{itemize}[leftmargin=1.35em,itemsep=1pt,topsep=1pt,parsep=0pt]
    \item $T$: source equation symbol.
    \item $\gamma$: source equation symbol.
    \item $u_{T}$: source equation symbol.
    \item $\mathcal{K}_{0007}$: canonical registry object for entry ID 0007.
    \item $\alpha_{T}$: source equation symbol.
    \item $Z$: source equation symbol.
    \item $P$: source equation symbol.
    \item $N$: normalization, count, or closure parameter in the entry relation.
    \item $R$: global radius, boundary radius, or topological scale.
    \item $A$: source equation symbol.
\end{itemize}
\endgroup
\end{minipage}\hfill
\begin{minipage}[t]{0.47\linewidth}
\tzpsection{Wolfram Form}
\begingroup
\scriptsize
\raggedright
\textbf{Source Wolfram model.}\par
\begin{Verbatim}[fontsize=\tiny,breaklines=true,breakanywhere=true]
(* TZPID ID0007 generated source Wolfram model *)
ClearAll[K0007, Cr0007, T, R, A, W, I, N];
eq0007$1 = HoldForm[alphaT == Derivative[1][dt]/dt];
eq0007$2 = HoldForm[alphaT^TZP == LimitGammap[Derivative[1][dt]/dt]];
eq0007$3 = HoldForm[alphaTNorm == NormalizeTRAWIN[Derivative[1][dt]/dt, LimitGammap[Derivative[1][dt]/dt], uT]];

K0007 = HoldForm[{eq0007$1, eq0007$2, eq0007$3}];

N = "normalized TRAWIN closure object for Temporal Displacement Factor";
I = "Normalized Temporal Operator integration/comparison layer";
W = "Approach-to-TZP Limit wave or operator carrier";
A = "Local-to-Global Time Ratio admissible action/amplitude condition";
R = "Approach-to-TZP Limit rotation/curl preservation layer";
T = "Local-to-Global Time Ratio local source condition";

Cr0007[expr_] := HoldForm[N[I[W[A[R[T[expr]]]]]]];

Result0007 = Cr0007[K0007];
\end{Verbatim}
\endgroup
\end{minipage}

\par\vspace{0.45\baselineskip}
\noindent\begin{minipage}[t]{\linewidth}
\begingroup
\small
\raggedright
\textbf{TRAWIN Operator Key.}\par
\noindent\textbf{Time/localization operator:} $\mathsf{T}\!\left[\mathrm{local}\right]$ Local-to-Global Time Ratio local source condition.\par
\noindent\textbf{Rotation/curl operator:} $\mathsf{R}\!\left[\nabla\times\right]$ Approach-to-TZP Limit rotation/curl preservation layer.\par
\noindent\textbf{Action/amplitude operator:} $\mathsf{A}\!\left[\mathrm{admissible}\right]$ Local-to-Global Time Ratio admissible action/amplitude condition.\par
\noindent\textbf{Wave/d'Alembertian operator:} $\mathsf{W}\!\left[\Box\right]$ Approach-to-TZP Limit wave or operator carrier.\par
\noindent\textbf{Integration operator:} $\mathsf{I}\!\left[\int\right]$ Normalized Temporal Operator integration/comparison layer.\par
\noindent\textbf{Normalization operator:} $\mathsf{N}\!\left[\mathrm{normalized}\right]$ normalized TRAWIN closure object for Temporal Displacement Factor.\par
\endgroup
\end{minipage}

\clearpage
\noindent\textbf{[ID 0007] Temporal Displacement Factor}\par
\vspace{0.35em}
\tzpsection{Derivation Layer}
\paragraph{Primary Equation.}
\begin{equation}
\alpha(t)=\frac{\mathrm{d}t'}{\mathrm{d}t}.
\end{equation}

\paragraph{Primary Derivation.}
Let $t$ denote ordinary parameter time and let $t'$ denote displaced or transformed time measured through the Trawinistic structure.

If $t'$ depends smoothly on $t$, then the infinitesimal rate of temporal displacement is given by the derivative:

\begin{equation}
\mathrm{d}t' = \alpha(t)\,\mathrm{d}t .
\end{equation}

Solving for $\alpha(t)$ gives:

\begin{equation}
\boxed{
\alpha(t)=\frac{\mathrm{d}t'}{\mathrm{d}t}.
}
\end{equation}

\paragraph{Interpretation.}
The temporal displacement factor measures how transformed time stretches, compresses, or shifts relative to ordinary time.

\par\vspace{0.35em}
\noindent\textbf{Derivable Consequences.}\quad
\begingroup
\footnotesize
\raggedright
The canonical equation anchors Temporal Displacement Factor by connecting the source condition to the derivation layer and proof certificate.
\par\endgroup

\tzpsection{Equation Support}
\noindent\textbf{1. Local-to-Global Time Ratio}\par
\[
\alpha_{T}
=
\frac{dt'}{dt}
\]
\noindent This defines the raw temporal comparison without yet committing to a singular limit.\par
\noindent\textbf{2. Approach-to-TZP Limit}\par
\[
\alpha_{T}^{\mathrm{TZP}}
=
\lim_{\gamma \to p} \frac{dt'}{dt}
\]
\noindent This localizes the factor to the one regime where metric degeneracy can distort the clock relationship.\par
\noindent\textbf{3. Normalized Temporal Operator}\par
\[
\alpha_{T}^{\mathrm{norm}}
=
\operatorname{Normalize}_{\mathrm{TRAWIN}}
\!\left(
\frac{dt'}{dt},\,
\lim_{\gamma \to p}\frac{dt'}{dt},\,
u_{T}
\right)
\]
\noindent This closes the progression by tying clock distortion to the framework's internal scale rather than a purely external parametrization.\par

\clearpage
\noindent\textbf{[ID 0007] Temporal Displacement Factor}\par
\vspace{0.35em}
\tzpsection{Dictionary}
\begingroup\small
\tzpsection{Definition}
Temporal Displacement Factor label id-007-temporal-displacement-factor is a registry definition in the active TZPID arc.

% BEGIN TZP CONTEXTUAL MEANING
\tzpsection{Contextual Meaning}
\begingroup\footnotesize\setlength{\emergencystretch}{3em}\sloppy
\textbf{Formal statement:} The temporal displacement factor AlphaT is defined by AlphaT = ratio dt dt , where t is the local time measured along a geodesic approaching p and t is the global time parameter. \textbf{Contextual meaning:} The entry reads Temporal Displacement Factor label id-007-temporal-displacement-factor as a controlled technical headword whose equation layer, proof route, and registry role are interpreted together. \textbf{Derivable consequences:} The entry preserves the named object as a reusable definition, connects it to the local proof route, and records the consequences implied by the equation chain. \textbf{Lemma support:} Local-to-Global Time Ratio; Approach-to-TZP Limit; Normalized Temporal Operator.\par
\endgroup
% END TZP CONTEXTUAL MEANING

\tzpsection{Technical Interpretation}
The temporal displacement factor measures how transformed time stretches, compresses, or shifts relative to ordinary time.

\tzpsection{Equation Grounding}
Equation anchors: alpha(t)=ratio d t d t , alpha(t)=ratio d t d t ,.

\tzpsection{Interpretive Role}
Once normalized, this factor becomes the clock correction that later Hamiltonians, actions, and propagation equations can use whenever they pass through or near the singular anchor.
\endgroup

\tzpsection{TZP Thesaurus}
\begingroup\footnotesize\sloppy
\begin{enumerate}[leftmargin=1.35em,itemsep=1pt,topsep=1pt,parsep=0pt]
\item \textbf{Time-Lag Displacement Factor Label Id}: entry-specific alternate wording for indexing, related literature, and local formal context.
\item \textbf{Time-Lag Time-Lag Formal Analogue}: entry-specific alternate wording for indexing, related literature, and local formal context.
\item \textbf{Time-Lag Terminology}: entry-specific alternate wording for indexing, related literature, and local formal context.
\end{enumerate}
\endgroup

\tzpsection{Encyclopedia Mathematica}
\begingroup\footnotesize\sloppy
The visual plate below is reserved for the source-specific equation and progression graphic for [ID 0007]. It is included only as a companion to the canonical equation, not as a substitute for the formal text.
\par\endgroup
\ifTZPDarkEdition
\tzpdictionaryvisualband{D:/TZPID/kdp_workflow/build/tikz_figures/ID0007_dictionary_figure_dark.pdf}
\fi

\clearpage
\begingroup\tzpsupportsetup
\noindent\textbf{\fontsize{10.8pt}{12.8pt}\selectfont [ID 0007] Constructive Logic and Proof Certificate}\par
\noindent\textbf{\fontsize{10pt}{12pt}\selectfont Temporal Displacement Factor}\par
\vspace{0.45em}
\begingroup
\footnotesize
\setlength{\parskip}{0.22em}
\setlength{\parindent}{0pt}
\noindent\textbf{Curry--Howard--Lambek Interpretation}\par
Under the Curry--Howard--Lambek correspondence, this registry entry is read as a typed proof
object: the stated source condition supplies the proposition, the derivation supplies the
morphism, and the proof certificate records the machine handles needed to check the obligation
in the formal registry.

\vspace{0.45em}
\noindent\textbf{CHL Proof}\par
\textbf{Statement:} the temporal displacement factor  is defined by  = dt'dt, where t' is the local
time measured along a geodesic approaching p and t is the...

\textbf{Logic Validation:}\\
\textbf{Proposition:} ``Temporal Displacement Factor is valid as a formal dynamical law when the Hamiltonian or
energy operator is well defined on the stated state space.''\\
\textbf{Proof:} The source statement identifies an energy, Hamiltonian, or vacuum-sector mechanism. The
CHL proof reads it as an operator claim: if the state space and localized terms are
defined, then the evolution or extraction channel is formally meaningful.

\textbf{VERDICT:} \ensuremath{\checkmark}\ \textbf{TRUE} --- formal dynamical-operator validation

\textbf{Formal Status:} \ensuremath{\checkmark}\ THEORETICAL -- source-backed CHL proof obligation

\textbf{Physical Status:} THEORETICAL -- requires model-specific empirical interpretation
\par\vspace{0.45em}
\noindent{\tiny\textbf{Isabelle/HOL.} \tzpplaincode{physics\_validity\_from\_model, external\_validity\_package, registry\_number\_matches, equation\_count\_matches}}\par
\ifTZPDarkEdition
\tzpproofvisualband{D:/TZPID/kdp_workflow/build/tikz_figures/ID0007_proof_certificate_figure_dark.pdf}
\fi
\endgroup
\endgroup
